% =====================================================================
% Resolução dos Desafios de Programação — Apêndice D
% "C: Como Programar" — Paul Deitel, Harvey Deitel, 6ª edição
% Programação de Jogos: Solução de Sudoku
% =====================================================================
\documentclass[11pt,a4paper,oneside]{report}

% --- Codificação e idioma ---
\usepackage[utf8]{inputenc}
\usepackage[T1]{fontenc}
\usepackage[brazil]{babel}
\usepackage{lmodern}

% --- Layout ---
\usepackage[a4paper,margin=2.4cm,headheight=22pt]{geometry}
\usepackage{setspace}
\onehalfspacing
\usepackage{microtype}
\usepackage{parskip}

% --- Cores, gráficos e boxes ---
\usepackage{xcolor}
\definecolor{deitelblue}{RGB}{0,82,147}
\definecolor{deitellightblue}{RGB}{210,232,250}
\definecolor{deitelgreen}{RGB}{34,120,57}
\definecolor{deitellightgreen}{RGB}{220,240,225}
\definecolor{deitelred}{RGB}{170,30,30}
\definecolor{deitellightred}{RGB}{250,225,225}
\definecolor{deitelorange}{RGB}{200,110,15}
\definecolor{deitellightorange}{RGB}{252,235,210}
\definecolor{deitelgray}{RGB}{75,75,75}
\definecolor{codebg}{RGB}{248,248,245}
\definecolor{codekw}{RGB}{0,82,147}
\definecolor{codecomment}{RGB}{0,128,0}
\definecolor{codestring}{RGB}{170,30,30}
\definecolor{codenum}{RGB}{120,120,120}

\usepackage[most]{tcolorbox}
\tcbuselibrary{listings,skins,breakable}

% --- Listings (estilo de código C no espírito Deitel) ---
\usepackage{listings}
\lstset{
    language=C,
    backgroundcolor=\color{codebg},
    basicstyle=\ttfamily\small,
    keywordstyle=\color{codekw}\bfseries,
    commentstyle=\color{codecomment}\itshape,
    stringstyle=\color{codestring},
    numberstyle=\tiny\color{codenum},
    numbers=left,
    numbersep=8pt,
    stepnumber=1,
    showstringspaces=false,
    breaklines=true,
    breakatwhitespace=true,
    tabsize=4,
    captionpos=b,
    frame=single,
    rulecolor=\color{deitelblue!50},
    framesep=4pt,
    xleftmargin=14pt,
    xrightmargin=4pt,
    aboveskip=10pt,
    belowskip=10pt,
    inputencoding=utf8,
    extendedchars=true,
    literate=
        {á}{{\'a}}1 {à}{{\`a}}1 {â}{{\^a}}1 {ã}{{\~a}}1
        {é}{{\'e}}1 {ê}{{\^e}}1 {í}{{\'i}}1
        {ó}{{\'o}}1 {ô}{{\^o}}1 {õ}{{\~o}}1
        {ú}{{\'u}}1 {ç}{{\c{c}}}1
        {Á}{{\'A}}1 {É}{{\'E}}1 {Í}{{\'I}}1
        {Ó}{{\'O}}1 {Ú}{{\'U}}1 {Ç}{{\c{C}}}1
}

% --- Caixas pedagógicas estilo Deitel ---
\newtcolorbox{boapratica}{
    enhanced, breakable,
    colback=deitellightgreen, colframe=deitelgreen,
    boxrule=0.6pt, arc=2pt,
    fonttitle=\bfseries\color{white},
    coltitle=white,
    title=\faIconHack{check} Boa Prática de Programação,
    fontupper=\small,
    left=8pt, right=8pt, top=4pt, bottom=4pt,
}
\newtcolorbox{erro}{
    enhanced, breakable,
    colback=deitellightred, colframe=deitelred,
    boxrule=0.6pt, arc=2pt,
    fonttitle=\bfseries\color{white},
    coltitle=white,
    title=\faIconHack{warning} Erro Comum de Programação,
    fontupper=\small,
    left=8pt, right=8pt, top=4pt, bottom=4pt,
}
\newtcolorbox{desempenho}{
    enhanced, breakable,
    colback=deitellightorange, colframe=deitelorange,
    boxrule=0.6pt, arc=2pt,
    fonttitle=\bfseries\color{white},
    coltitle=white,
    title=\faIconHack{bolt} Dica de Desempenho,
    fontupper=\small,
    left=8pt, right=8pt, top=4pt, bottom=4pt,
}
\newtcolorbox{engenharia}{
    enhanced, breakable,
    colback=deitellightblue, colframe=deitelblue,
    boxrule=0.6pt, arc=2pt,
    fonttitle=\bfseries\color{white},
    coltitle=white,
    title=\faIconHack{gear} Observação de Engenharia de Software,
    fontupper=\small,
    left=8pt, right=8pt, top=4pt, bottom=4pt,
}
\newtcolorbox{enunciado}{
    enhanced, breakable,
    colback=white, colframe=deitelblue,
    boxrule=1pt, arc=3pt,
    fonttitle=\bfseries\color{white},
    coltitle=white,
    title=Enunciado do desafio,
    fontupper=\small,
    left=8pt, right=8pt, top=4pt, bottom=4pt,
}

% Implementação simples de "ícones" via texto entre colchetes (sem font-awesome para evitar dependência)
\newcommand{\faIconHack}[1]{[\textsc{#1}]}

% --- Capítulos e seções ---
\usepackage{titlesec}
\titleformat{\chapter}[display]
    {\normalfont\huge\bfseries\color{deitelblue}}
    {\filright\Large APÊNDICE \thechapter}
    {18pt}
    {\Huge}
\titleformat{\section}
    {\normalfont\Large\bfseries\color{deitelblue}}
    {\thesection}{0.8em}{}
\titleformat{\subsection}
    {\normalfont\large\bfseries\color{deitelgray}}
    {\thesubsection}{0.7em}{}

\renewcommand{\thechapter}{D}
\renewcommand{\thesection}{D.\arabic{section}}

% --- Cabeçalhos e rodapés ---
\usepackage{fancyhdr}
\pagestyle{fancy}
\fancyhf{}
\fancyhead[L]{\small\color{deitelgray}\textit{Resolução — Apêndice D}}
\fancyhead[R]{\small\color{deitelgray}\textit{C: Como Programar — Deitel \& Deitel, 6ª ed.}}
\fancyfoot[C]{\small\color{deitelgray}\thepage}
\renewcommand{\headrulewidth}{0.4pt}
\renewcommand{\headrule}{\hbox to\headwidth{\color{deitelblue}\leaders\hrule height \headrulewidth\hfill}}

% --- Hyperref por último ---
\usepackage[colorlinks=true,
    linkcolor=deitelblue,
    citecolor=deitelblue,
    urlcolor=deitelblue,
    pdftitle={Resolução do Apêndice D - C Como Programar - Deitel},
    pdfauthor={Resolução em estilo Deitel},
    pdfsubject={Programação de Jogos: Solução de Sudoku em C}
]{hyperref}

% --- Comandos auxiliares ---
\newcommand{\codigo}[1]{\texttt{\color{deitelblue}#1}}
\newcommand{\funcao}[1]{\texttt{\color{deitelblue}\textbf{#1}}}

\begin{document}

% =====================================================================
% CAPA
% =====================================================================
\begin{titlepage}
\centering
\vspace*{1.5cm}
{\Huge\bfseries\color{deitelblue} Resolução Comentada}\\[0.5cm]
{\LARGE\bfseries\color{deitelblue} Apêndice D}\\[0.4cm]
{\Large\color{deitelgray} Programação de Jogos:}\\[0.2cm]
{\Large\color{deitelgray} Solução de Sudoku}\\[1.5cm]
\rule{0.7\textwidth}{1pt}\\[0.8cm]
{\large\textit{Baseado na obra}}\\[0.4cm]
{\Large\bfseries C: Como Programar}\\[0.2cm]
{\large Paul Deitel \quad$\bullet$\quad Harvey Deitel}\\[0.2cm]
{\normalsize 6\textsuperscript{a} edição --- Pearson Prentice Hall, 2011}\\[1.2cm]
\rule{0.7\textwidth}{1pt}\\[1.5cm]

\begin{tcolorbox}[
    colback=deitellightblue,
    colframe=deitelblue,
    width=0.85\textwidth,
    arc=4pt,
    boxrule=1pt
]
\small
\textbf{Sobre esta resolução.} Este documento apresenta a resolução completa
e didática dos oito desafios de programação propostos no Apêndice D da obra
\textit{C: Como Programar}. Cada desafio é desenvolvido no estilo
incremental característico dos autores --- enunciado, análise, algoritmo,
implementação comentada em C, saída esperada e observações pedagógicas
(boas práticas, erros comuns, dicas de desempenho). Como o Apêndice D
sintetiza conceitos que vão de arrays bidimensionais (Cap. 6) até recursão
e ponteiros (Caps. 5 e 7), o material aqui apresentado utiliza apenas
recursos da linguagem C apresentados até esses capítulos.
\end{tcolorbox}

\vfill
{\small\color{deitelgray}\today}
\end{titlepage}

% =====================================================================
% SUMÁRIO
% =====================================================================
\tableofcontents
\thispagestyle{fancy}
\newpage

% =====================================================================
% APÊNDICE D
% =====================================================================
\chapter{Programação de Jogos: Solução de Sudoku}

\section{Visão geral e infraestrutura comum}

O Apêndice~D apresenta um conjunto de desafios de programação cujo objetivo
é, ao mesmo tempo, divertir o leitor e exercitar todos os mecanismos que
foram introduzidos nos capítulos anteriores: arrays bidimensionais, funções,
ponteiros, estruturas de iteração aninhadas, geração de números pseudoaleatórios
e --- onde o problema exigir --- recursão. Antes de atacarmos cada desafio
individualmente, é prudente construir uma pequena infraestrutura comum a
todos eles. Essa estratégia é, aliás, a mesma adotada em todo o livro:
identificamos as funcionalidades reutilizáveis, programamo-las uma única vez
e, a partir daí, compomos soluções cada vez mais sofisticadas.

\begin{engenharia}
Reaproveitar funções utilitárias é a essência da boa engenharia de software.
Em vez de duplicar a verificação ``este movimento é válido?'' em cada
solucionador, escrevemos a função uma única vez e a invocamos onde for
necessário. Quando uma falha aparecer (ou um aprimoramento for desejado),
basta corrigir um único ponto do código.
\end{engenharia}

\subsection{Convenções de representação do tabuleiro}

Seguindo a recomendação do Apêndice~D, representamos o tabuleiro como um
array bidimensional de inteiros de dimensões $10 \times 10$, ignorando a
\emph{linha~0} e a \emph{coluna~0}. Assim, a célula da $i$-ésima linha e
$j$-ésima coluna corresponde a \codigo{s[i][j]} com $1 \le i,j \le 9$,
em conformidade com a numeração tradicional da comunidade do Sudoku. Uma
célula vazia é representada pelo valor \codigo{0}; uma célula preenchida
contém um dos dígitos \codigo{1} a \codigo{9}.

\begin{boapratica}
Definir constantes simbólicas (\textit{macros}) para o tamanho do tabuleiro
elimina ``números mágicos'' espalhados pelo código. Caso desejássemos, no
futuro, generalizar para um tabuleiro $16 \times 16$, alteraríamos apenas
as macros.
\end{boapratica}

\subsection{Função utilitária \texttt{ehMovimentoValido}}

Praticamente todos os solucionadores precisarão responder à pergunta:
``\emph{Se eu colocar o dígito $d$ na célula \codigo{s[i][j]}, isso violará
alguma regra do Sudoku?}'' Em vez de reverificar o tabuleiro inteiro a cada
tentativa, é muito mais eficiente verificar apenas a linha~$i$, a coluna~$j$
e a grade de $3 \times 3$ à qual \codigo{s[i][j]} pertence. A
Listagem~\ref{lst:infra} apresenta o cabeçalho do programa, as constantes
e essa função auxiliar.

\begin{lstlisting}[caption={Macros, cabeçalhos e a função \texttt{ehMovimentoValido}.},label={lst:infra}]
/* sudoku_infra.h -- infraestrutura comum aos desafios do Apêndice D. */
#ifndef SUDOKU_INFRA_H
#define SUDOKU_INFRA_H

#include <stdio.h>
#include <stdlib.h>
#include <time.h>

#define LADO       9   /* dimensão do Sudoku padrão */
#define DIM        10  /* dimensão do array; ignoramos índice 0 */
#define VAZIO      0   /* valor que indica célula em branco */
#define BLOCO      3   /* dimensão das sub-grades 3x3 */

/* Verifica se colocar o dígito d em s[linha][col] mantém o tabuleiro */
/* consistente com as regras do Sudoku. Retorna 1 se sim, 0 se não.   */
int ehMovimentoValido(int s[DIM][DIM], int linha, int col, int d)
{
    int k, i, j;
    int inicioL, inicioC;

    /* 1) o dígito já apareceu na linha? */
    for (k = 1; k <= LADO; k++)
        if (s[linha][k] == d) return 0;

    /* 2) o dígito já apareceu na coluna? */
    for (k = 1; k <= LADO; k++)
        if (s[k][col] == d) return 0;

    /* 3) o dígito já apareceu na grade 3x3 que contém [linha][col]? */
    inicioL = ((linha - 1) / BLOCO) * BLOCO + 1;
    inicioC = ((col   - 1) / BLOCO) * BLOCO + 1;
    for (i = inicioL; i < inicioL + BLOCO; i++)
        for (j = inicioC; j < inicioC + BLOCO; j++)
            if (s[i][j] == d) return 0;

    return 1;   /* nenhum conflito */
}
#endif
\end{lstlisting}

A aritmética para localizar a grade $3 \times 3$ merece atenção. A expressão
\codigo{((linha - 1) / 3) * 3 + 1} faz exatamente o que precisamos: subtraímos
$1$ para que as linhas $1, 2, 3$ caiam em $0, 1, 2$; dividimos por $3$ (divisão
inteira) para obter o índice do bloco ($0$, $1$ ou $2$); multiplicamos por $3$
e somamos $1$ para voltar à numeração de $1$ a $9$, obtendo a primeira linha
da grade ($1$, $4$ ou $7$, respectivamente). O mesmo raciocínio se aplica à
coluna.

\begin{erro}
Um erro frequente é escrever \codigo{(linha / 3) * 3 + 1} sem o ajuste
\codigo{- 1}. Para \codigo{linha == 3}, isso produziria $4$, quando o início
correto do bloco é $1$. Sempre teste a aritmética de bloco com os valores
de fronteira $1$, $3$, $4$, $6$, $7$ e $9$.
\end{erro}

\subsection{Função utilitária \texttt{imprimirTabuleiro}}

Como o padrão da linguagem C não inclui recursos gráficos, optamos por uma
representação textual cuidadosa, com separadores que destacam as sub-grades
de $3 \times 3$. A Listagem~\ref{lst:print} apresenta a implementação.

\begin{lstlisting}[caption={Impressão formatada do tabuleiro.},label={lst:print}]
/* Imprime o tabuleiro com separadores entre as sub-grades 3x3. */
void imprimirTabuleiro(int s[DIM][DIM])
{
    int i, j;

    printf("\n    1 2 3 | 4 5 6 | 7 8 9\n");
    printf("  +-------+-------+-------+\n");
    for (i = 1; i <= LADO; i++) {
        printf("%d | ", i);
        for (j = 1; j <= LADO; j++) {
            if (s[i][j] == VAZIO)
                printf(". ");
            else
                printf("%d ", s[i][j]);
            if (j % 3 == 0) printf("| ");
        }
        printf("\n");
        if (i % 3 == 0)
            printf("  +-------+-------+-------+\n");
    }
}
\end{lstlisting}

Com essa pequena infraestrutura --- macros de domínio, função de validação
incremental e função de impressão --- estamos prontos para enfrentar cada
um dos desafios. Nas seções a seguir, supomos que essas funções já estão
disponíveis (em um cabeçalho \codigo{sudoku\_infra.h} incluído em cada
programa).

% =====================================================================
\section{Desafio 1: a função \texttt{validSudoku}}

\begin{enunciado}
Crie a função \texttt{int validSudoku( int sudokuBoard[10][10] )} que recebe
um tabuleiro de Sudoku como um array bidimensional de inteiros (ignorando a
linha 0 e a coluna 0). A função deve retornar \texttt{1} se um tabuleiro
\emph{completo} for válido, \texttt{2} se um tabuleiro \emph{parcialmente
completo} for válido e \texttt{0} se o tabuleiro for \emph{inválido}.
\end{enunciado}

\subsection{Análise do problema}

Antes de escrever uma única linha de código, é essencial entender com
precisão o que ``válido'' significa neste contexto. Um Sudoku é válido
quando, ao mesmo tempo:

\begin{enumerate}
    \item \textbf{Todas as células contêm valores em $\{0, 1, 2, \ldots, 9\}$.}
        O valor $0$ representa célula vazia; qualquer outro valor fora desse
        intervalo é, por si só, um erro.
    \item \textbf{Em cada linha, cada dígito de 1 a 9 aparece, no máximo, uma vez.}
        Em um tabuleiro completo, todos os nove dígitos devem aparecer
        exatamente uma vez; em um tabuleiro parcial, basta que não haja
        repetição entre as células já preenchidas.
    \item \textbf{O mesmo vale para cada coluna.}
    \item \textbf{O mesmo vale para cada uma das nove grades de $3 \times 3$.}
\end{enumerate}

A distinção entre ``completo e válido'' (retorno \codigo{1}), ``parcial e
válido'' (retorno \codigo{2}) e ``inválido'' (retorno \codigo{0}) sugere
naturalmente o seguinte plano: percorremos todas as células para detectar
duplicatas; ao mesmo tempo, registramos se existe pelo menos uma célula
vazia. Ao final, se algum conflito foi detectado, retornamos \codigo{0};
se nenhum conflito foi detectado e existem células vazias, retornamos
\codigo{2}; se nenhum conflito foi detectado e todas as 81 células estão
preenchidas, retornamos \codigo{1}.

\subsection{Algoritmo}

\begin{enumerate}
    \item \textbf{Validar o domínio dos valores.} Percorra todas as células.
        Se alguma célula contiver valor fora de $\{0, 1, \ldots, 9\}$,
        retorne \codigo{0}. Caso contrário, marque \codigo{temVazio = 1}
        se encontrar pelo menos uma célula com valor \codigo{0}.
    \item \textbf{Verificar repetições nas nove linhas.} Para cada linha,
        utilize um array auxiliar \codigo{presente[10]} (índices 1 a 9)
        para marcar quais dígitos já apareceram. Se um dígito já estiver
        marcado, retorne \codigo{0}.
    \item \textbf{Verificar repetições nas nove colunas.} Procedimento
        análogo ao das linhas, trocando os papéis de linha e coluna.
    \item \textbf{Verificar repetições nas nove sub-grades $3 \times 3$.}
        Use dois laços externos (sobre os blocos $0,1,2$) e dois internos
        (sobre as três linhas e três colunas dentro do bloco).
    \item \textbf{Determinar o veredito.} Se chegamos até aqui, não há
        conflito. Retorne \codigo{2} se \codigo{temVazio == 1}, ou
        \codigo{1} se o tabuleiro estiver completo.
\end{enumerate}

\begin{boapratica}
Reinicializar o vetor \codigo{presente} a cada linha (ou coluna, ou bloco)
é o passo mais facilmente esquecido nessa rotina. Se o array não for zerado,
um dígito que apareceu na linha~1 pode ser interpretado como repetição na
linha~2, gerando falsos positivos.
\end{boapratica}

\subsection{Implementação}

\begin{lstlisting}[caption={Implementação completa da função \texttt{validSudoku}.},label={lst:valid}]
/* validSudoku.c -- desafio 1 do Apêndice D.                       */
/* Retorna: 1 se completo e válido, 2 se parcial e válido,         */
/*          0 se houver qualquer violação das regras.              */
#include "sudoku_infra.h"

int validSudoku(int sudokuBoard[DIM][DIM])
{
    int i, j, k, valor;
    int presente[DIM];   /* presente[d] != 0 se d já foi visto */
    int temVazio = 0;
    int blocoLinha, blocoCol;

    /* 1) Validação de domínio dos valores. */
    for (i = 1; i <= LADO; i++) {
        for (j = 1; j <= LADO; j++) {
            valor = sudokuBoard[i][j];
            if (valor < 0 || valor > LADO)
                return 0;             /* valor inválido */
            if (valor == VAZIO)
                temVazio = 1;
        }
    }

    /* 2) Verificação das nove linhas. */
    for (i = 1; i <= LADO; i++) {
        for (k = 0; k < DIM; k++) presente[k] = 0;
        for (j = 1; j <= LADO; j++) {
            valor = sudokuBoard[i][j];
            if (valor != VAZIO) {
                if (presente[valor])  /* já apareceu */
                    return 0;
                presente[valor] = 1;
            }
        }
    }

    /* 3) Verificação das nove colunas. */
    for (j = 1; j <= LADO; j++) {
        for (k = 0; k < DIM; k++) presente[k] = 0;
        for (i = 1; i <= LADO; i++) {
            valor = sudokuBoard[i][j];
            if (valor != VAZIO) {
                if (presente[valor])
                    return 0;
                presente[valor] = 1;
            }
        }
    }

    /* 4) Verificação das nove sub-grades 3x3. */
    for (blocoLinha = 0; blocoLinha < BLOCO; blocoLinha++) {
        for (blocoCol = 0; blocoCol < BLOCO; blocoCol++) {
            for (k = 0; k < DIM; k++) presente[k] = 0;
            for (i = 1; i <= BLOCO; i++) {
                for (j = 1; j <= BLOCO; j++) {
                    valor = sudokuBoard[blocoLinha*BLOCO + i]
                                       [blocoCol  *BLOCO + j];
                    if (valor != VAZIO) {
                        if (presente[valor])
                            return 0;
                        presente[valor] = 1;
                    }
                }
            }
        }
    }

    /* 5) Veredito final. */
    return temVazio ? 2 : 1;
}
\end{lstlisting}

\subsection{Programa de teste e saída esperada}

A Listagem~\ref{lst:valid_test} apresenta um pequeno programa que constrói
três tabuleiros --- um inválido, um parcial e um completo --- e invoca
\codigo{validSudoku} em cada um deles.

\begin{lstlisting}[caption={Programa de teste para \texttt{validSudoku}.},label={lst:valid_test}]
int main(void)
{
    /* Tabuleiro completo e VÁLIDO. */
    int completo[DIM][DIM] = {
        { 0 },
        { 0, 5,3,4, 6,7,8, 9,1,2 },
        { 0, 6,7,2, 1,9,5, 3,4,8 },
        { 0, 1,9,8, 3,4,2, 5,6,7 },
        { 0, 8,5,9, 7,6,1, 4,2,3 },
        { 0, 4,2,6, 8,5,3, 7,9,1 },
        { 0, 7,1,3, 9,2,4, 8,5,6 },
        { 0, 9,6,1, 5,3,7, 2,8,4 },
        { 0, 2,8,7, 4,1,9, 6,3,5 },
        { 0, 3,4,5, 2,8,6, 1,7,9 }
    };
    /* Tabuleiro parcial e válido (várias células vazias). */
    int parcial[DIM][DIM] = {
        { 0 },
        { 0, 5,3,0, 0,7,0, 0,0,0 },
        { 0, 6,0,0, 1,9,5, 0,0,0 },
        { 0, 0,9,8, 0,0,0, 0,6,0 },
        { 0, 8,0,0, 0,6,0, 0,0,3 },
        { 0, 4,0,0, 8,0,3, 0,0,1 },
        { 0, 7,0,0, 0,2,0, 0,0,6 },
        { 0, 0,6,0, 0,0,0, 2,8,0 },
        { 0, 0,0,0, 4,1,9, 0,0,5 },
        { 0, 0,0,0, 0,8,0, 0,7,9 }
    };
    /* Tabuleiro INVÁLIDO: dois 5 na primeira linha. */
    int invalido[DIM][DIM] = {
        { 0 },
        { 0, 5,5,0, 0,7,0, 0,0,0 },  /* erro proposital */
        { 0, 6,0,0, 1,9,5, 0,0,0 },
        { 0, 0,9,8, 0,0,0, 0,6,0 },
        { 0, 8,0,0, 0,6,0, 0,0,3 },
        { 0, 4,0,0, 8,0,3, 0,0,1 },
        { 0, 7,0,0, 0,2,0, 0,0,6 },
        { 0, 0,6,0, 0,0,0, 2,8,0 },
        { 0, 0,0,0, 4,1,9, 0,0,5 },
        { 0, 0,0,0, 0,8,0, 0,7,9 }
    };

    printf("Tabuleiro completo  -> validSudoku = %d\n",
           validSudoku(completo));
    printf("Tabuleiro parcial   -> validSudoku = %d\n",
           validSudoku(parcial));
    printf("Tabuleiro invalido  -> validSudoku = %d\n",
           validSudoku(invalido));

    return 0;
}
\end{lstlisting}

\textbf{Saída esperada:}
\begin{lstlisting}[language=,basicstyle=\ttfamily\small]
Tabuleiro completo  -> validSudoku = 1
Tabuleiro parcial   -> validSudoku = 2
Tabuleiro invalido  -> validSudoku = 0
\end{lstlisting}

\begin{desempenho}
A complexidade dessa função é $\Theta(n^2)$, com $n = 9$ --- ou seja, um
custo constante na prática. Como ela será invocada milhões de vezes pelos
solucionadores por força bruta, ainda assim vale a pena evitar trabalho
desnecessário. Em situações em que apenas \emph{uma} célula foi alterada,
a função \codigo{ehMovimentoValido} (apresentada na Seção D.1) faz a mesma
verificação examinando apenas $27$ células, e por isso é claramente
preferível dentro de laços de busca.
\end{desempenho}

% =====================================================================
\section{Desafio 2: solucionador por força bruta com retrocesso}

\begin{enunciado}
A versão mais inteligente da técnica de força bruta consiste em verificar
cada dígito antes de colocá-lo: se o tabuleiro permanece válido, prossiga
para a próxima célula; se não, experimente outro dígito; se nenhum dos nove
dígitos funcionar, retorne à célula anterior e troque seu valor. Implemente
um solucionador com essa estratégia.
\end{enunciado}

\subsection{Análise do problema}

A descrição acima é a definição clássica do algoritmo de \emph{retrocesso}
(em inglês, \textit{backtracking}). A ideia é simples e poderosa: tentamos
construir uma solução um passo de cada vez; sempre que percebemos que o
caminho atual é inviável, ``voltamos atrás'' e tentamos outro. Como cada
tentativa é validada antes de prosseguirmos, podamos enormes porções do
espaço de busca --- $9^{81}$ possibilidades cegas tornam-se algumas dezenas
de milhares de tentativas inteligentes em um Sudoku típico.

A recursão é, sem dúvida, a forma mais elegante de expressar esse algoritmo.
Aprendemos no Capítulo~5 que uma função recursiva resolve um problema
maior reduzindo-o a uma versão menor de si mesmo. Aqui o ``problema menor''
é o tabuleiro com \emph{uma célula a mais} preenchida.

\begin{engenharia}
O retrocesso é onipresente em problemas combinatórios: o Passeio do Cavalo,
as Oito Rainhas, palavras cruzadas, geração de permutações --- todos seguem
o mesmo esqueleto: ``tente, valide, recurse, desfaça''. Reconhecer esse
padrão é o primeiro passo para resolver dezenas de problemas aparentemente
distintos com o mesmo arcabouço mental.
\end{engenharia}

\subsection{Algoritmo}

\begin{enumerate}
    \item \textbf{Encontre a primeira célula vazia.} Percorra o tabuleiro
        nas direções natural (linhas crescentes e, dentro de cada linha,
        colunas crescentes).
    \item \textbf{Caso-base.} Se não houver célula vazia, o tabuleiro está
        completo e válido (pois cada dígito foi validado ao ser colocado).
        Retorne sucesso.
    \item \textbf{Passo recursivo.} Para cada dígito $d$ de $1$ a $9$:
        \begin{enumerate}
            \item Use \codigo{ehMovimentoValido} para verificar se $d$ pode
                  ser colocado na célula corrente.
            \item Em caso afirmativo, atribua \codigo{s[i][j] = d} e chame
                  recursivamente o solucionador.
            \item Se a chamada recursiva retornar sucesso, propague sucesso.
            \item Caso contrário, \textbf{desfaça} a atribuição
                  (\codigo{s[i][j] = VAZIO}) e tente o próximo dígito.
        \end{enumerate}
    \item Se nenhum dígito funcionou, retorne falha (isso obriga a chamada
        anterior a tentar outro dígito).
\end{enumerate}

\begin{erro}
Esquecer de \emph{desfazer} a atribuição (\codigo{s[i][j] = VAZIO}) antes
de tentar o próximo dígito é o erro mais comum em algoritmos de retrocesso.
Sem essa restauração, o tabuleiro acumula valores ``fantasmas'' das
tentativas anteriores, e o algoritmo simplesmente falha em todos os
caminhos subsequentes.
\end{erro}

\subsection{Implementação}

\begin{lstlisting}[caption={Solucionador por força bruta com retrocesso.},label={lst:bt}]
/* sudoku_backtracking.c -- desafio 2 do Apêndice D.            */
/* Resolve o tabuleiro modificando-o in place. Retorna 1 se uma */
/* solução foi encontrada, 0 caso contrário.                     */
#include "sudoku_infra.h"

int resolverBacktracking(int s[DIM][DIM])
{
    int i, j, d;

    /* 1) Localiza a primeira célula vazia em ordem linha-coluna. */
    for (i = 1; i <= LADO; i++) {
        for (j = 1; j <= LADO; j++) {
            if (s[i][j] == VAZIO) {

                /* 2) Tenta cada dígito de 1 a 9. */
                for (d = 1; d <= LADO; d++) {
                    if (ehMovimentoValido(s, i, j, d)) {
                        s[i][j] = d;                /* tente   */
                        if (resolverBacktracking(s))/* recurse */
                            return 1;
                        s[i][j] = VAZIO;            /* desfaça */
                    }
                }
                /* Nenhum dígito serviu: avisa o nível anterior. */
                return 0;
            }
        }
    }
    /* Nenhuma célula vazia: tabuleiro completo (caso-base). */
    return 1;
}

int main(void)
{
    /* Quebra-cabeça clássico publicado pelo The Times (Londres). */
    int s[DIM][DIM] = {
        { 0 },
        { 0, 5,3,0, 0,7,0, 0,0,0 },
        { 0, 6,0,0, 1,9,5, 0,0,0 },
        { 0, 0,9,8, 0,0,0, 0,6,0 },
        { 0, 8,0,0, 0,6,0, 0,0,3 },
        { 0, 4,0,0, 8,0,3, 0,0,1 },
        { 0, 7,0,0, 0,2,0, 0,0,6 },
        { 0, 0,6,0, 0,0,0, 2,8,0 },
        { 0, 0,0,0, 4,1,9, 0,0,5 },
        { 0, 0,0,0, 0,8,0, 0,7,9 }
    };

    printf("Tabuleiro inicial:");
    imprimirTabuleiro(s);

    if (resolverBacktracking(s)) {
        printf("\nSolucao encontrada:");
        imprimirTabuleiro(s);
    } else {
        printf("\nNao ha solucao para esse tabuleiro.\n");
    }
    return 0;
}
\end{lstlisting}

\subsection{Saída esperada}

\begin{lstlisting}[language=,basicstyle=\ttfamily\footnotesize]
Tabuleiro inicial:
    1 2 3 | 4 5 6 | 7 8 9
  +-------+-------+-------+
1 | 5 3 . | . 7 . | . . . |
2 | 6 . . | 1 9 5 | . . . |
3 | . 9 8 | . . . | . 6 . |
  +-------+-------+-------+
4 | 8 . . | . 6 . | . . 3 |
5 | 4 . . | 8 . 3 | . . 1 |
6 | 7 . . | . 2 . | . . 6 |
  +-------+-------+-------+
7 | . 6 . | . . . | 2 8 . |
8 | . . . | 4 1 9 | . . 5 |
9 | . . . | . 8 . | . 7 9 |
  +-------+-------+-------+

Solucao encontrada:
    1 2 3 | 4 5 6 | 7 8 9
  +-------+-------+-------+
1 | 5 3 4 | 6 7 8 | 9 1 2 |
2 | 6 7 2 | 1 9 5 | 3 4 8 |
3 | 1 9 8 | 3 4 2 | 5 6 7 |
  +-------+-------+-------+
4 | 8 5 9 | 7 6 1 | 4 2 3 |
5 | 4 2 6 | 8 5 3 | 7 9 1 |
6 | 7 1 3 | 9 2 4 | 8 5 6 |
  +-------+-------+-------+
7 | 9 6 1 | 5 3 7 | 2 8 4 |
8 | 2 8 7 | 4 1 9 | 6 3 5 |
9 | 3 4 5 | 2 8 6 | 1 7 9 |
  +-------+-------+-------+
\end{lstlisting}

\begin{desempenho}
Este solucionador resolve, em frações de segundo, virtualmente qualquer
Sudoku publicado em jornais. Sudokus considerados ``mais difíceis do mundo''
podem exigir centenas de milhares de chamadas recursivas, mas mesmo nesses
casos o tempo total é tipicamente inferior a um segundo em um computador
moderno. A simplicidade desse algoritmo --- uma única função de cerca de
20 linhas --- contrasta com as estratégias heurísticas elaboradas, e
ilustra um princípio importante: \emph{quando o hardware é farto, a
simplicidade vence}.
\end{desempenho}

\subsection{Comparação com a versão ``ingênua'' de 81 laços aninhados}

O Apêndice~D menciona uma técnica alternativa: aninhar 81 laços \codigo{for},
cada um varrendo de $1$ a $9$. O número total de combinações é $9^{81}$,
ou aproximadamente $1{,}97 \times 10^{77}$. Para se ter ideia da escala,
mesmo a $10^{12}$ tentativas por segundo --- desempenho de um
supercomputador --- testar todas elas levaria cerca de $6 \times 10^{57}$
anos, muitas ordens de magnitude superior à idade estimada do universo.

A versão com retrocesso ``poda'' esse espaço de busca explorando apenas
caminhos viáveis. Em termos práticos, o número de combinações efetivamente
visitadas cai para a faixa de algumas dezenas de milhares. Esse é o poder
de uma decisão arquitetural correta.

% =====================================================================
\section{Desafio 3: gerador por permutações aleatórias}

\begin{enunciado}
Escreva a função \texttt{void permutations(int sudokuBoard[10][10])} que
recebe um array bidimensional de $10 \times 10$ e preenche cada uma das
nove linhas da grade $9 \times 9$ com uma permutação selecionada
aleatoriamente dos dígitos de $1$ a $9$. Em seguida, use \texttt{validSudoku}
em laço até obter um Sudoku válido. Inclua uma opção em \texttt{validSudoku}
que a faça verificar apenas colunas e sub-grades de $3 \times 3$.
\end{enunciado}

\subsection{Análise do problema}

Esta abordagem tem um forte sabor de \emph{Monte Carlo}: tentamos
configurações aleatórias até que uma delas, por sorte, seja válida. A
probabilidade de uma única tentativa produzir um Sudoku válido é
extremamente baixa, mas \emph{não nula}; com tempo suficiente, o gerador
eventualmente terá sucesso.

A ideia, descrita no apêndice, é gerar uma permutação dos dígitos de $1$
a $9$ para cada linha. Como cada linha é, por construção, uma permutação,
não há repetições \emph{dentro de uma mesma linha}. A função
\codigo{validSudoku} só precisa, portanto, verificar duas coisas: que não
há repetições nas colunas e que não há repetições nas sub-grades $3 \times 3$.
Daí a opção adicional que vamos introduzir.

\begin{desempenho}
A probabilidade de uma combinação aleatória de nove permutações constituir
um Sudoku válido é da ordem de $10^{-21}$. Ou seja, na prática, este
gerador é didático --- ilustra perfeitamente o que ``força bruta'' significa
--- mas é absolutamente impraticável como gerador real. Em programas sérios,
combina-se este método com retrocesso (vide Seção D.8).
\end{desempenho}

\subsection{Geração de uma permutação aleatória}

O Apêndice~D descreve o método ingênuo: ``escolha um dígito; para o próximo,
escolha repetidamente até que seja diferente do anterior; e assim por
diante''. Implementamos exatamente esse método na Listagem~\ref{lst:perm}.

\begin{lstlisting}[caption={Geração de uma linha como permutação aleatória.},label={lst:perm}]
/* Preenche linha[1..9] com uma permutação aleatória de 1..9.   */
void permutaLinha(int linha[DIM])
{
    int usado[DIM];
    int pos, d;

    /* Marca todos os dígitos como ainda não utilizados. */
    for (d = 0; d < DIM; d++) usado[d] = 0;

    /* Para cada posição 1..9, sorteia repetidamente até obter  */
    /* um dígito ainda não utilizado.                           */
    for (pos = 1; pos <= LADO; pos++) {
        do {
            d = (rand() % LADO) + 1;   /* sorteia em [1, 9] */
        } while (usado[d]);
        usado[d] = 1;
        linha[pos] = d;
    }
}

/* Aplica permutaLinha a cada uma das nove linhas do tabuleiro. */
void permutations(int sudokuBoard[DIM][DIM])
{
    int i;
    for (i = 1; i <= LADO; i++)
        permutaLinha(sudokuBoard[i]);
}
\end{lstlisting}

\begin{boapratica}
O laço \codigo{do...while} aqui é justo: \emph{precisamos} executar o sorteio
pelo menos uma vez antes de testar a condição. Esse é o critério clássico
para escolher entre \codigo{while} e \codigo{do...while} --- se o teste
deve ocorrer \emph{depois} da primeira execução do corpo, use
\codigo{do...while}.
\end{boapratica}

\subsection{Variante de \texttt{validSudoku} com verificação seletiva}

O apêndice nos pede para introduzir uma opção que faça a função verificar
\emph{apenas} colunas e sub-grades. Há várias maneiras de implementar isso;
optamos por uma flag inteira que recebe um padrão de bits (estilo clássico
de C): bit~0 ativa a verificação de linhas, bit~1 a de colunas, bit~2 a de
sub-grades. Para verificar todas as três, passamos \codigo{0x7}; para
verificar apenas colunas e sub-grades, passamos \codigo{0x6}.

\begin{lstlisting}[caption={\texttt{validSudoku} com flags de verificação seletiva.},label={lst:valid_flags}]
#define VERIFICA_LINHAS  0x1
#define VERIFICA_COLUNAS 0x2
#define VERIFICA_BLOCOS  0x4
#define VERIFICA_TUDO    (VERIFICA_LINHAS|VERIFICA_COLUNAS|VERIFICA_BLOCOS)

int validSudokuSel(int s[DIM][DIM], int opcoes)
{
    int i, j, k, valor;
    int presente[DIM];
    int temVazio = 0;
    int bl, bc;

    /* Domínio dos valores (sempre verificado). */
    for (i = 1; i <= LADO; i++)
        for (j = 1; j <= LADO; j++) {
            valor = s[i][j];
            if (valor < 0 || valor > LADO) return 0;
            if (valor == VAZIO) temVazio = 1;
        }

    if (opcoes & VERIFICA_LINHAS)
        for (i = 1; i <= LADO; i++) {
            for (k = 0; k < DIM; k++) presente[k] = 0;
            for (j = 1; j <= LADO; j++) {
                valor = s[i][j];
                if (valor != VAZIO) {
                    if (presente[valor]) return 0;
                    presente[valor] = 1;
                }
            }
        }

    if (opcoes & VERIFICA_COLUNAS)
        for (j = 1; j <= LADO; j++) {
            for (k = 0; k < DIM; k++) presente[k] = 0;
            for (i = 1; i <= LADO; i++) {
                valor = s[i][j];
                if (valor != VAZIO) {
                    if (presente[valor]) return 0;
                    presente[valor] = 1;
                }
            }
        }

    if (opcoes & VERIFICA_BLOCOS)
        for (bl = 0; bl < BLOCO; bl++)
            for (bc = 0; bc < BLOCO; bc++) {
                for (k = 0; k < DIM; k++) presente[k] = 0;
                for (i = 1; i <= BLOCO; i++)
                    for (j = 1; j <= BLOCO; j++) {
                        valor = s[bl*BLOCO+i][bc*BLOCO+j];
                        if (valor != VAZIO) {
                            if (presente[valor]) return 0;
                            presente[valor] = 1;
                        }
                    }
            }
    return temVazio ? 2 : 1;
}
\end{lstlisting}

\subsection{Programa gerador completo}

\begin{lstlisting}[caption={Gerador por permutações aleatórias --- prova de conceito.},label={lst:gerador_perm}]
int main(void)
{
    int s[DIM][DIM] = {{0}};
    long tentativas = 0;
    const long LIMITE = 5000000L;  /* limite prático */

    srand((unsigned) time(NULL));

    do {
        permutations(s);
        tentativas++;
    } while (validSudokuSel(s, VERIFICA_COLUNAS | VERIFICA_BLOCOS) != 1
             && tentativas < LIMITE);

    if (tentativas < LIMITE) {
        printf("Sudoku gerado apos %ld tentativas:\n", tentativas);
        imprimirTabuleiro(s);
    } else {
        printf("Limite de %ld tentativas atingido sem sucesso.\n", LIMITE);
        printf("(Isso e esperado: a probabilidade e da ordem de 10^-21.)\n");
    }
    return 0;
}
\end{lstlisting}

\subsection{Discussão do experimento}

Executar este programa é instrutivo: na esmagadora maioria das execuções,
o limite de tentativas é atingido sem sucesso. Esse comportamento, longe
de ser uma falha do código, é uma \emph{característica fundamental} do
método: por isso é chamado de ``força bruta''. O número aproximado de
arranjos cuja primeira linha é qualquer permutação fixa e cujas demais
linhas constituem um Sudoku válido é, segundo cálculos publicados na
literatura, $6{,}67 \times 10^{21}$. Como existem $9!^9 \approx 1{,}1 \times 10^{50}$
combinações de nove permutações independentes, a probabilidade de ``acerto
ao acaso'' é da ordem de $6 \times 10^{-29}$ --- isto é, virtualmente zero.

\begin{engenharia}
Algoritmos puros de Monte Carlo, sem nenhuma poda, são úteis apenas como
\emph{baseline} pedagógico. O próximo passo natural é combinar o sorteio
aleatório com algum mecanismo de poda --- exatamente o que faremos quando
unirmos esta ideia ao retrocesso, no gerador de quebra-cabeças (Seção D.8).
\end{engenharia}

\subsection{Variante prática --- usando retrocesso para construir um Sudoku válido}

Como exercício adicional, observe que podemos construir Sudokus válidos
de modo \emph{determinístico} embaralhando a ordem em que o solucionador
do Desafio~2 testa os dígitos. Em vez de tentar sempre $1, 2, \ldots, 9$,
geramos uma permutação aleatória da sequência $\{1, 2, \ldots, 9\}$ para
cada célula vazia. O resultado é um Sudoku completo e \emph{aleatório},
gerado em milissegundos. Esta é a abordagem realmente útil na prática;
a função \codigo{permutations} pura, com retorno via \codigo{validSudoku},
serve para demonstrar o conceito de força bruta.

% =====================================================================
\section{Desafio 4: solucionador ``fácil'' --- estratégia dos singletons}

\begin{enunciado}
Implemente um solucionador para Sudokus ``fáceis'' baseado na repetição
das estratégias humanas: para cada célula vazia, mantenha a lista de
dígitos ainda possíveis (com base nos valores já fixados na linha, coluna
e sub-grade); identifique singletons (célula com uma única possibilidade
restante) e singletons ocultos (dígito que só cabe em uma célula da unidade)
e os fixe; repita até que nenhuma estratégia produza progresso.
\end{enunciado}

\subsection{Análise do problema}

As estratégias descritas no apêndice --- singletons, duplas, duplas ocultas,
triplas, triplas ocultas --- podem ser implementadas em qualquer ordem,
porque qualquer uma delas, ao agir, simplifica o tabuleiro para as demais.
Vamos concentrar nossa atenção nas duas mais produtivas: o singleton
``nu'' (célula com uma única possibilidade) e o singleton ``oculto''
(dígito que aparece como possível em apenas uma célula de uma linha, coluna
ou sub-grade). Essas duas estratégias resolvem, sozinhas, a grande maioria
dos Sudokus de níveis fácil e médio.

\textbf{Representação das possibilidades.} Para cada célula vazia, queremos
um conjunto de dígitos. A representação clássica em C é uma \emph{máscara
de bits}: usamos um \codigo{int} em que o bit $(d - 1)$ está ligado se, e
somente se, o dígito $d$ ainda é possível para aquela célula. Isso permite
operações conjuntivas elegantes:
\begin{itemize}
    \item \codigo{(mask \& (1 << (d-1)))} testa se $d$ pertence ao conjunto;
    \item \codigo{mask |= (1 << (d-1))} adiciona $d$ ao conjunto;
    \item \codigo{mask \&= \~{}(1 << (d-1))} remove $d$ do conjunto;
    \item contar quantos bits estão ligados conta o tamanho do conjunto.
\end{itemize}

\begin{engenharia}
Máscaras de bits são onipresentes em C quando precisamos manipular
``conjuntos de até 32 (ou 64) elementos''. O custo de memória é mínimo,
as operações são triviais e tudo cabe em registradores do processador.
Em situações didáticas pode-se usar um array de 9 inteiros booleanos por
célula, mas o ganho de clareza paga-se com perda de desempenho. Para
Sudoku, a máscara é a escolha natural.
\end{engenharia}

\subsection{Algoritmo}

\begin{enumerate}
    \item \textbf{Calcule as possibilidades.} Para cada célula vazia,
        examine sua linha, sua coluna e sua sub-grade; o conjunto dos
        dígitos ausentes desses três conjuntos é a lista de possibilidades.
    \item \textbf{Aplique singletons nus.} Se a lista de alguma célula
        contém um único dígito, fixe esse dígito na célula.
    \item \textbf{Aplique singletons ocultos.} Para cada linha, cada coluna
        e cada sub-grade, e para cada dígito $d$ de $1$ a $9$, conte em
        quantas células vazias dessa unidade o dígito $d$ é possível.
        Se a contagem for $1$, fixe $d$ naquela célula.
    \item \textbf{Repita.} Após qualquer alteração no tabuleiro, recalcule
        as possibilidades e reaplique as estratégias. Pare quando uma
        iteração completa não produzir nenhuma alteração.
    \item \textbf{Veredito.} Se o tabuleiro está completo, sucesso. Caso
        contrário, este Sudoku exige técnicas mais avançadas (vide
        Seções D.5 e D.6).
\end{enumerate}

\subsection{Implementação}

\begin{lstlisting}[caption={Funções auxiliares para manipulação de bitmasks.},label={lst:bitmask}]
/* Retorna o número de bits ligados em mask. */
int contarBits(int mask)
{
    int n = 0;
    while (mask) {
        n += mask & 1;
        mask >>= 1;
    }
    return n;
}

/* Se mask tem exatamente um bit ligado, retorna o dígito       */
/* correspondente (1..9); caso contrário, retorna 0.            */
int unicoDigito(int mask)
{
    int d;
    if (contarBits(mask) != 1) return 0;
    for (d = 1; d <= LADO; d++)
        if (mask & (1 << (d - 1))) return d;
    return 0;
}
\end{lstlisting}

\begin{lstlisting}[caption={Cálculo da matriz de possibilidades para cada célula.},label={lst:poss}]
/* Para cada célula vazia, registra em possib[i][j] uma máscara */
/* de bits dos dígitos ainda permitidos.                        */
void calcularPossibilidades(int s[DIM][DIM], int possib[DIM][DIM])
{
    int i, j, d;
    for (i = 1; i <= LADO; i++)
        for (j = 1; j <= LADO; j++) {
            if (s[i][j] != VAZIO) {
                possib[i][j] = 0;        /* já preenchida */
                continue;
            }
            possib[i][j] = 0;
            for (d = 1; d <= LADO; d++)
                if (ehMovimentoValido(s, i, j, d))
                    possib[i][j] |= (1 << (d - 1));
        }
}
\end{lstlisting}

\begin{lstlisting}[caption={Aplicação dos singletons nus.},label={lst:singleton_nu}]
/* Percorre o tabuleiro fixando toda célula que tenha exatamente */
/* uma possibilidade. Retorna 1 se houve mudança, 0 se não.       */
int aplicarSingletonsNus(int s[DIM][DIM], int possib[DIM][DIM])
{
    int i, j, d, mudou = 0;
    for (i = 1; i <= LADO; i++)
        for (j = 1; j <= LADO; j++)
            if (s[i][j] == VAZIO) {
                d = unicoDigito(possib[i][j]);
                if (d != 0) {
                    s[i][j] = d;
                    possib[i][j] = 0;
                    mudou = 1;
                }
            }
    return mudou;
}
\end{lstlisting}

\begin{lstlisting}[caption={Aplicação dos singletons ocultos.},label={lst:singleton_oculto}]
/* Para cada linha, coluna ou bloco, procura um dígito que só    */
/* pode ser colocado em uma única célula da unidade --- e o fixa.*/
int aplicarSingletonsOcultos(int s[DIM][DIM], int possib[DIM][DIM])
{
    int i, j, d, mudou = 0;
    int contagem, ii, jj, posI = 0, posJ = 0, br, bc;

    /* Singletons ocultos em LINHAS. */
    for (i = 1; i <= LADO; i++)
        for (d = 1; d <= LADO; d++) {
            contagem = 0;
            for (j = 1; j <= LADO; j++)
                if (s[i][j] == VAZIO
                    && (possib[i][j] & (1 << (d - 1)))) {
                    contagem++;
                    posJ = j;
                }
            if (contagem == 1) {
                s[i][posJ] = d;
                possib[i][posJ] = 0;
                mudou = 1;
            }
        }

    /* Singletons ocultos em COLUNAS. */
    for (j = 1; j <= LADO; j++)
        for (d = 1; d <= LADO; d++) {
            contagem = 0;
            for (i = 1; i <= LADO; i++)
                if (s[i][j] == VAZIO
                    && (possib[i][j] & (1 << (d - 1)))) {
                    contagem++;
                    posI = i;
                }
            if (contagem == 1) {
                s[posI][j] = d;
                possib[posI][j] = 0;
                mudou = 1;
            }
        }

    /* Singletons ocultos em SUB-GRADES 3x3. */
    for (br = 0; br < BLOCO; br++)
        for (bc = 0; bc < BLOCO; bc++)
            for (d = 1; d <= LADO; d++) {
                contagem = 0;
                for (ii = 1; ii <= BLOCO; ii++)
                    for (jj = 1; jj <= BLOCO; jj++) {
                        i = br * BLOCO + ii;
                        j = bc * BLOCO + jj;
                        if (s[i][j] == VAZIO
                            && (possib[i][j] & (1 << (d - 1)))) {
                            contagem++;
                            posI = i; posJ = j;
                        }
                    }
                if (contagem == 1) {
                    s[posI][posJ] = d;
                    possib[posI][posJ] = 0;
                    mudou = 1;
                }
            }
    return mudou;
}
\end{lstlisting}

\begin{lstlisting}[caption={Solucionador completo para Sudokus fáceis.},label={lst:solver_facil}]
/* Aplica as estratégias repetidamente, recalculando             */
/* possibilidades após cada rodada de alterações.                */
/* Retorna 1 se o tabuleiro foi totalmente resolvido.            */
int resolverFacil(int s[DIM][DIM])
{
    int possib[DIM][DIM];
    int progresso, i, j, vazias;

    do {
        calcularPossibilidades(s, possib);
        progresso = aplicarSingletonsNus(s, possib);

        /* Recalcule antes da próxima estratégia para refletir   */
        /* as células recém-fixadas.                              */
        calcularPossibilidades(s, possib);
        progresso |= aplicarSingletonsOcultos(s, possib);
    } while (progresso);

    /* Conta células ainda vazias. */
    vazias = 0;
    for (i = 1; i <= LADO; i++)
        for (j = 1; j <= LADO; j++)
            if (s[i][j] == VAZIO) vazias++;

    return vazias == 0;
}
\end{lstlisting}

\subsection{Discussão e limitações}

Esta versão resolve, sem qualquer suposição ou retrocesso, virtualmente
todos os Sudokus rotulados como ``fácil'' por jornais e revistas. Para
Sudokus médios, ela costuma ficar muito próxima da solução --- frequentemente
deixa apenas duas ou três células sem decisão. Para Sudokus difíceis,
porém, ela trava muito antes do fim.

\begin{boapratica}
Recalcular as possibilidades \emph{antes} de cada estratégia evita que
informações estagnadas levem o algoritmo a deduções inválidas. O custo é
pequeno (algumas dezenas de operações), e a robustez aumenta enormemente.
\end{boapratica}

\subsection{Comentário sobre duplas, triplas e suas versões ocultas}

As estratégias de duplas, triplas, duplas ocultas e triplas ocultas
descritas no Apêndice~D são extensões naturais do raciocínio dos
singletons: em vez de procurar ``uma célula com uma possibilidade'',
procuramos ``duas células que dividem exatamente duas possibilidades'',
ou ``dois dígitos que só cabem em duas células de uma unidade''. A lógica
de implementação é mais elaborada, mas sempre se reduz a percorrer cada
unidade (linha, coluna ou sub-grade) e, dentro dela, examinar combinações
de tamanho $2$ ou $3$. Para os propósitos pedagógicos deste apêndice, os
singletons (nus e ocultos) capturam a essência da técnica; deixamos
duplas e triplas como extensão recomendada ao leitor.

% =====================================================================
\section{Desafio 5: heurística ``esperar antes de tomar uma decisão''}

\begin{enunciado}
Implemente a heurística do tipo \emph{esperar antes de tomar uma decisão}:
$(i)$ associe a cada célula vazia uma lista de possibilidades; $(ii)$
caracterize o estado do tabuleiro contando o total de posicionamentos
ainda possíveis; $(iii)$ para cada movimento candidato, avalie a contagem
do tabuleiro \emph{após} esse movimento; $(iv)$ escolha o movimento que
deixar a contagem mais alta (em caso de empate, sorteie).
\end{enunciado}

\subsection{Análise do problema}

A heurística é exatamente a oposta da intuição comum. Quando o problema é
puramente combinatório --- como no Passeio do Cavalo do Capítulo~6 e nas
Oito Rainhas --- a tentação é avançar rapidamente fechando opções. A
filosofia ``mantenha suas opções abertas'' diz o contrário: \emph{prefira
a jogada que preserve o maior número de futuros possíveis}. A justificativa
é simples: cada jogada que ``fecha'' o tabuleiro aumenta a probabilidade
de que, mais adiante, alguma célula fique sem nenhum dígito viável ---
um beco sem saída.

A medida quantitativa proposta é a soma, sobre todas as células vazias do
tabuleiro, do número de dígitos ainda permitidos. Essa contagem dá um
``orçamento de liberdade'' para o tabuleiro. Entre dois movimentos
candidatos, preferimos o que deixar a contagem \emph{maior}.

\begin{engenharia}
Heurísticas, por natureza, não oferecem garantias. Esta, em particular,
não é sempre superior ao retrocesso ingênuo --- mas é especialmente eficaz
em Sudokus de dificuldade intermediária, onde reduz dramaticamente o
número de jogadas hipotéticas. É também uma excelente ferramenta de
estudo: ela ilustra como uma decisão simples (qual jogada fazer agora?)
pode ser informada por uma medida global do estado do problema.
\end{engenharia}

\subsection{Algoritmo}

\begin{enumerate}
    \item Defina \codigo{contarPossibilidades(s)} que soma, para cada
        célula vazia, o número de dígitos $d \in \{1, \ldots, 9\}$
        admissíveis (via \codigo{ehMovimentoValido}).
    \item Para cada par (célula vazia, dígito admissível):
        \begin{enumerate}
            \item Atribua hipoteticamente o dígito à célula.
            \item Compute a contagem global resultante.
            \item Restaure a célula a vazio.
        \end{enumerate}
    \item Selecione o movimento de maior contagem. Em caso de empate,
        sorteie um dos empatados.
    \item Efetue o movimento escolhido (definitivamente).
    \item Repita até que o tabuleiro esteja completo ou até que algum
        movimento leve a contagem a zero (beco sem saída).
    \item Em becos sem saída, retroceda ao último ponto de decisão ---
        ou seja, combine a heurística com retrocesso.
\end{enumerate}

\subsection{Implementação}

\begin{lstlisting}[caption={Contagem global de possibilidades.},label={lst:cnt}]
/* Soma, sobre cada célula vazia, o número de dígitos válidos.   */
int contarPossibilidades(int s[DIM][DIM])
{
    int i, j, d, total = 0;
    for (i = 1; i <= LADO; i++)
        for (j = 1; j <= LADO; j++)
            if (s[i][j] == VAZIO)
                for (d = 1; d <= LADO; d++)
                    if (ehMovimentoValido(s, i, j, d))
                        total++;
    return total;
}
\end{lstlisting}

\begin{lstlisting}[caption={Heurística ``esperar antes de tomar uma decisão''.},label={lst:wait}]
/* Estrutura para representar um movimento candidato. */
typedef struct {
    int i, j, d, contagem;
} Movimento;

/* Tenta cada movimento possível e elege o que preserva a maior */
/* contagem global de possibilidades. Retorna 1 se conseguiu    */
/* efetuar um movimento; 0 se não há mais movimentos válidos.    */
int escolherMovimentoEspera(int s[DIM][DIM])
{
    Movimento empate[729];     /* no máximo 81 * 9 movimentos */
    int nEmpate = 0;
    int melhor = -1;
    int i, j, d, c;

    for (i = 1; i <= LADO; i++)
        for (j = 1; j <= LADO; j++)
            if (s[i][j] == VAZIO)
                for (d = 1; d <= LADO; d++)
                    if (ehMovimentoValido(s, i, j, d)) {
                        s[i][j] = d;
                        c = contarPossibilidades(s);
                        s[i][j] = VAZIO;

                        if (c > melhor) {
                            melhor = c;
                            nEmpate = 0;
                            empate[nEmpate].i = i;
                            empate[nEmpate].j = j;
                            empate[nEmpate].d = d;
                            empate[nEmpate].contagem = c;
                            nEmpate++;
                        } else if (c == melhor) {
                            empate[nEmpate].i = i;
                            empate[nEmpate].j = j;
                            empate[nEmpate].d = d;
                            empate[nEmpate].contagem = c;
                            nEmpate++;
                        }
                    }

    if (nEmpate == 0)
        return 0;   /* nenhum movimento válido */

    /* Em caso de empate, sorteia. */
    int idx = rand() % nEmpate;
    s[ empate[idx].i ][ empate[idx].j ] = empate[idx].d;
    return 1;
}

int resolverHeuristica(int s[DIM][DIM])
{
    int i, j, vazias;

    do {
        if (!escolherMovimentoEspera(s))
            return 0;          /* beco sem saída */

        vazias = 0;
        for (i = 1; i <= LADO; i++)
            for (j = 1; j <= LADO; j++)
                if (s[i][j] == VAZIO) vazias++;
    } while (vazias > 0);

    return 1;
}
\end{lstlisting}

\subsection{Discussão de complexidade}

A cada movimento, examinamos no máximo $81 \times 9 = 729$ candidatos e,
para cada um, recomputamos uma contagem que custa também $81 \times 9$
operações. O custo de uma única decisão é, portanto, da ordem de
$729 \times 729 \approx 5{,}3 \times 10^5$ operações. Como cada Sudoku
exige até $81$ decisões, o custo total fica em torno de $4{,}3 \times 10^7$
operações no pior caso --- ainda perfeitamente factível em um computador
moderno.

\begin{desempenho}
O laço de contagem é o ponto quente do algoritmo. Em produção, vale a
pena manter uma matriz de possibilidades (vide Desafio~4) e atualizá-la
incrementalmente: quando colocamos $d$ em \codigo{s[i][j]}, basta remover
o bit correspondente das máscaras das células da mesma linha, coluna e
bloco. Isso reduz o custo por movimento a $O(27)$ em vez de $O(729)$.
\end{desempenho}

\begin{erro}
Esquecer de restaurar a célula a \codigo{VAZIO} após a avaliação hipotética
é um bug particularmente difícil de diagnosticar: o tabuleiro vai
acumulando movimentos fantasmas, e o algoritmo passa a ``ver'' contagens
sistematicamente mais baixas, levando a escolhas erradas. Use sempre o
padrão ``\textbf{tente, meça, desfaça}''.
\end{erro}

\subsection{Combinando heurística e retrocesso}

A heurística pura, sem retrocesso, é capaz de resolver Sudokus difíceis
em muitas execuções, mas eventualmente atinge becos sem saída. A forma
robusta é manter uma pilha de pontos de decisão: cada vez que escolhemos
entre vários candidatos empatados, salvamos o estado do tabuleiro. Se
mais adiante chegamos a um beco sem saída, restauramos o estado salvo
mais recente e tentamos outro empatado. Em essência, transformamos a
heurística em um guia para o retrocesso: ela diz \emph{qual ramo} explorar
primeiro, enquanto o retrocesso garante a corretude.

% =====================================================================
\section{Desafio 6: heurística da antecipação (look-ahead)}

\begin{enunciado}
A heurística da antecipação é um aprimoramento da heurística de
``esperar antes de tomar uma decisão''. No caso de empate, antecipe
mais uma colocação: para cada movimento empatado, avalie qual seria
a contagem do tabuleiro após \emph{duas} jogadas e prefira aquele cujo
melhor desdobramento for o mais alto.
\end{enunciado}

\subsection{Análise do problema}

A heurística da Seção D.6 às vezes não consegue distinguir entre vários
movimentos com o mesmo número de possibilidades imediatas. Quando isso
ocorre, fazer um sorteio é honesto, mas ingênuo --- estamos ignorando
informação que está apenas ``um passo mais à frente''. A heurística da
antecipação propõe: para cada candidato empatado, simule a próxima
jogada (a melhor possível, segundo o mesmo critério) e use a contagem
\emph{após dois movimentos} como critério de desempate.

\begin{engenharia}
Estamos aplicando, em pequena escala, um princípio universal de
inteligência artificial: \emph{busca em profundidade limitada}. O xadrez,
o Go e dezenas de outros jogos usam a mesma ideia, geralmente em
profundidades muito maiores e com poda (alfa-beta) para evitar explorar
ramos óbvios. A diferença é que, no Sudoku, a função de avaliação é
simples e exata (a contagem global), enquanto em jogos adversariais ela
é apenas uma estimativa.
\end{engenharia}

\subsection{Algoritmo}

\begin{enumerate}
    \item Compute todos os movimentos candidatos como antes, registrando
        sua contagem imediata.
    \item Identifique a melhor contagem imediata.
    \item Filtre apenas os movimentos empatados nesse máximo.
    \item Para cada empatado, simule a jogada e, no tabuleiro resultante,
        compute a melhor contagem imediata da próxima jogada (sem
        efetivamente jogá-la --- apenas avalie). Esse valor é a
        \emph{contagem antecipada} do candidato.
    \item Restaure a célula simulada.
    \item Selecione o candidato com a maior contagem antecipada. Em caso
        de novo empate, sorteie.
\end{enumerate}

\subsection{Implementação}

\begin{lstlisting}[caption={Heurística com antecipação de uma jogada.},label={lst:lookahead}]
/* Dado um tabuleiro, encontra a melhor contagem imediata        */
/* alcançável por algum movimento (sem efetuá-lo).                */
int melhorContagemImediata(int s[DIM][DIM])
{
    int i, j, d, c, melhor = -1;
    for (i = 1; i <= LADO; i++)
        for (j = 1; j <= LADO; j++)
            if (s[i][j] == VAZIO)
                for (d = 1; d <= LADO; d++)
                    if (ehMovimentoValido(s, i, j, d)) {
                        s[i][j] = d;
                        c = contarPossibilidades(s);
                        s[i][j] = VAZIO;
                        if (c > melhor) melhor = c;
                    }
    return melhor;
}

/* Substitui escolherMovimentoEspera com critério de antecipação. */
int escolherMovimentoAntecipacao(int s[DIM][DIM])
{
    Movimento empate[729];
    int nEmpate = 0;
    int melhorImediata = -1;
    int i, j, d, c;

    /* Fase 1: coletar empatados na melhor contagem imediata.    */
    for (i = 1; i <= LADO; i++)
        for (j = 1; j <= LADO; j++)
            if (s[i][j] == VAZIO)
                for (d = 1; d <= LADO; d++)
                    if (ehMovimentoValido(s, i, j, d)) {
                        s[i][j] = d;
                        c = contarPossibilidades(s);
                        s[i][j] = VAZIO;

                        if (c > melhorImediata) {
                            melhorImediata = c;
                            nEmpate = 0;
                        }
                        if (c == melhorImediata) {
                            empate[nEmpate].i = i;
                            empate[nEmpate].j = j;
                            empate[nEmpate].d = d;
                            empate[nEmpate].contagem = c;
                            nEmpate++;
                        }
                    }

    if (nEmpate == 0)
        return 0;

    /* Fase 2: se houver um único empatado, jogue-o.             */
    if (nEmpate == 1) {
        s[empate[0].i][empate[0].j] = empate[0].d;
        return 1;
    }

    /* Fase 3: desempate por antecipação.                         */
    Movimento finalistas[729];
    int nFinalistas = 0;
    int melhorAntecipada = -1;

    for (int k = 0; k < nEmpate; k++) {
        s[empate[k].i][empate[k].j] = empate[k].d;
        int antecipada = melhorContagemImediata(s);
        s[empate[k].i][empate[k].j] = VAZIO;

        if (antecipada > melhorAntecipada) {
            melhorAntecipada = antecipada;
            nFinalistas = 0;
        }
        if (antecipada == melhorAntecipada) {
            finalistas[nFinalistas++] = empate[k];
        }
    }

    int idx = rand() % nFinalistas;
    s[finalistas[idx].i][finalistas[idx].j] = finalistas[idx].d;
    return 1;
}
\end{lstlisting}

\subsection{Custo e benefício}

A antecipação custa caro: para cada empatado, examinamos novamente
todos os movimentos candidatos. No pior caso, isso eleva o custo por
decisão a $O(n^4)$, com $n = 9$. Na prática, porém, o número de empatados
costuma ser pequeno (de $2$ a $5$ na maioria dos tabuleiros), de modo que
o tempo extra é tolerável. O benefício é uma redução notável no número
de becos sem saída: tabuleiros que exigiriam retrocesso resolvem-se,
muitas vezes, sem qualquer back-up.

\begin{boapratica}
Heurísticas com look-ahead profundo são, em geral, sujeitas à lei dos
rendimentos decrescentes: olhar dois movimentos à frente compensa quase
sempre; olhar três compensa às vezes; olhar quatro raramente paga seu
custo. Comece simples; aprofunde apenas se medições indicarem benefício.
\end{boapratica}

% =====================================================================
\section{Desafio 7: criando quebra-cabeças com células vazias}

\begin{enunciado}
Com um Sudoku completo em mãos, esvazie algumas células para criar um
quebra-cabeça. Considere duas estratégias: $(i)$ escolha aleatória de
células a esvaziar; $(ii)$ esvaziamento simétrico, usando a célula
``reflexiva'' --- definida como $s[10 - i][10 - j]$ para a célula
$s[i][j]$ --- de modo a deixar o tabuleiro com simetria visual.
\end{enunciado}

\subsection{Análise do problema}

A geração de Sudokus completos é o que fizemos no Desafio~2 (com
retrocesso) ou no Desafio~3 (com permutações aleatórias). A criação de
um \emph{quebra-cabeça} é o passo seguinte: removemos algumas células,
deixando para o jogador o trabalho de redescobri-las. O número de
células removidas é um \emph{proxy} da dificuldade: tabuleiros mais
``vazios'' tendem a ser mais difíceis, embora a relação não seja estrita
(há Sudokus com poucos espaços que são surpreendentemente difíceis e
vice-versa).

\textbf{Simetria reflexiva.} Sudokus publicados em jornais quase sempre
apresentam um padrão simétrico de células fixas, por estética. A célula
``reflexiva'' de $s[i][j]$ é $s[10 - i][10 - j]$ --- esta é a
involução $i \mapsto 10 - i$, $j \mapsto 10 - j$, equivalente a uma
rotação de $180^\circ$ em torno do centro do tabuleiro. Note dois casos
importantes:
\begin{itemize}
    \item Se $i = 5$ e $j = 5$, a célula é seu próprio reflexo (centro).
    \item Se $(i, j)$ e $(10 - i, 10 - j)$ são distintas, esvaziar uma
        obriga a esvaziar a outra para preservar a simetria.
\end{itemize}

\subsection{Algoritmo (esvaziamento aleatório simples)}

\begin{enumerate}
    \item Receba um Sudoku completo e uma quantidade-alvo
        \codigo{numCelulasParaEsvaziar}.
    \item Enquanto o número de células esvaziadas for menor que o alvo:
    \begin{enumerate}
        \item Sorteie aleatoriamente $i, j \in \{1, \ldots, 9\}$.
        \item Se \codigo{s[i][j]} ainda estiver preenchida, esvazie-a.
    \end{enumerate}
\end{enumerate}

\subsection{Algoritmo (esvaziamento simétrico)}

\begin{enumerate}
    \item Sorteie $i, j \in \{1, \ldots, 9\}$.
    \item Calcule $i' = 10 - i$ e $j' = 10 - j$.
    \item Se ambas \codigo{s[i][j]} e \codigo{s[i'][j']} estão preenchidas:
        esvazie ambas. Trate o caso especial $i = i'$ e $j = j'$ (centro)
        como um único esvaziamento.
\end{enumerate}

\subsection{Implementação}

\begin{lstlisting}[caption={Dois esquemas de esvaziamento de células.},label={lst:gerador_puzzle}]
/* Esvazia aleatoriamente n células de um Sudoku completo.       */
void esvaziarAleatorio(int s[DIM][DIM], int n)
{
    int esvaziadas = 0;
    int i, j;
    while (esvaziadas < n) {
        i = (rand() % LADO) + 1;
        j = (rand() % LADO) + 1;
        if (s[i][j] != VAZIO) {
            s[i][j] = VAZIO;
            esvaziadas++;
        }
    }
}

/* Esvazia n células de modo a preservar a simetria reflexiva.   */
/* n é o número-alvo de células esvaziadas; o algoritmo pode      */
/* parar com até n+1 esvaziadas (quando o último par é central). */
void esvaziarSimetrico(int s[DIM][DIM], int n)
{
    int esvaziadas = 0;
    int i, j, i2, j2;
    while (esvaziadas < n) {
        i = (rand() % LADO) + 1;
        j = (rand() % LADO) + 1;
        i2 = (LADO + 1) - i;   /* 10 - i */
        j2 = (LADO + 1) - j;   /* 10 - j */

        if (s[i][j] != VAZIO && s[i2][j2] != VAZIO) {
            if (i == i2 && j == j2) {
                /* Célula central: conta como um único esvaziamento. */
                s[i][j] = VAZIO;
                esvaziadas++;
            } else {
                s[i][j]   = VAZIO;
                s[i2][j2] = VAZIO;
                esvaziadas += 2;
            }
        }
    }
}
\end{lstlisting}

\begin{erro}
A inicialização do gerador de números aleatórios com \codigo{srand} é
frequentemente esquecida. Sem ela, \codigo{rand()} produz a mesma
sequência em cada execução --- e o programa gera sempre o mesmo
``Sudoku aleatório''. Use \codigo{srand( (unsigned) time(NULL) )} logo
no início de \codigo{main}.
\end{erro}

\subsection{Sobre dificuldade e solvibilidade}

Um detalhe sutil: o esvaziamento desimpedido pode gerar tabuleiros com
\emph{mais de uma} solução. Por exemplo, se removermos todas as células
de uma sub-grade, ela ficará vazia, e quaisquer duas atribuições
permutáveis das nove células gerarão Sudokus distintos. Para gerar
quebra-cabeças ``apropriados'', o ideal é, após cada esvaziamento,
verificar se o tabuleiro ainda tem solução única (vide Seção D.9); se
não, restaurar a última remoção. Esse aprimoramento transforma o gerador
em uma ferramenta verdadeiramente útil.

% =====================================================================
\section{Desafio 8: demonstrando que um Sudoku tem solução única}

\begin{enunciado}
Os Sudokus publicados em geral têm exatamente uma solução. Desenvolva um
meio de demonstrar que um jogo específico tem exatamente uma solução.
\end{enunciado}

\subsection{Análise do problema}

A primeira tentação é tentar gerar \emph{todas} as soluções e contar.
Para Sudokus mal definidos com poucas pistas, isso pode levar tempo
proibitivo --- existem Sudokus mal definidos com mais de $10^{15}$
soluções. Felizmente, a pergunta que precisamos responder é mais
modesta: ``há exatamente uma solução?''. Isso significa que basta
descobrir se há \emph{pelo menos duas}; ao encontrar a segunda, podemos
parar imediatamente.

A estratégia é adaptar o solucionador por retrocesso do Desafio~2: em
vez de retornar logo após encontrar uma solução, ele continua a busca
e conta. Como só nos interessa a distinção ``uma solução / mais de uma'',
abortamos a busca assim que a contagem atinge $2$. Essa poda é o que
torna o algoritmo prático.

\begin{engenharia}
Este é um exemplo clássico de \emph{poda baseada no objetivo}: não
precisamos da resposta exata para uma pergunta caríssima; bastaria
distinguir entre ``zero'', ``um'' e ``dois ou mais''. Reformular a
pergunta para que esse limiar seja pequeno transforma um problema
intratável em um problema rotineiro.
\end{engenharia}

\subsection{Algoritmo}

\begin{enumerate}
    \item Receba o tabuleiro a verificar.
    \item Copie-o (a verificação não deve alterar o original).
    \item Execute uma busca por retrocesso na cópia, contando soluções.
    \item Aborte a busca assim que o contador atinja $2$.
    \item Retorne ``única'' se o contador for $1$; ``múltiplas'' se for
        $2$; ``sem solução'' se for $0$.
\end{enumerate}

\subsection{Implementação}

\begin{lstlisting}[caption={Contagem limitada de soluções e teste de unicidade.},label={lst:unique}]
/* Conta soluções até o limite indicado. Modifica s (caller deve copiar). */
int contarSolucoes(int s[DIM][DIM], int limite)
{
    int i, j, d, total = 0;
    for (i = 1; i <= LADO; i++) {
        for (j = 1; j <= LADO; j++) {
            if (s[i][j] == VAZIO) {
                for (d = 1; d <= LADO; d++)
                    if (ehMovimentoValido(s, i, j, d)) {
                        s[i][j] = d;
                        total += contarSolucoes(s, limite - total);
                        s[i][j] = VAZIO;
                        if (total >= limite)
                            return total;     /* poda */
                    }
                return total;
            }
        }
    }
    /* Sem células vazias: encontramos uma solução completa. */
    return 1;
}

/* Verifica se o tabuleiro tem exatamente uma solução.           */
/* Retorna:                                                       */
/*   0 -> sem solução                                             */
/*   1 -> exatamente uma solução                                  */
/*   2 -> duas ou mais soluções                                   */
int contarUnicidade(int original[DIM][DIM])
{
    int copia[DIM][DIM];
    int i, j;

    /* Copia o tabuleiro para preservar o original. */
    for (i = 0; i < DIM; i++)
        for (j = 0; j < DIM; j++)
            copia[i][j] = original[i][j];

    return contarSolucoes(copia, 2);
}
\end{lstlisting}

\subsection{Discussão da poda}

A chave dessa solução está na linha \codigo{if (total >= limite) return total}.
Note que ``$\codigo{limite} = 2$'' significa: ``conte até no máximo duas
soluções; assim que achar a segunda, pare''. A passagem
\codigo{limite - total} na chamada recursiva propaga o ``orçamento de
soluções'' restante. Isso transforma o algoritmo, que no pior caso é
exponencial, em algo cujo custo escala apenas com a árvore necessária
para encontrar até duas soluções.

\begin{boapratica}
Sempre copiar a entrada antes de executar uma operação destrutiva
permite que o chamador continue a usar o tabuleiro original. A regra
geral é: ``funções de consulta não devem modificar seus argumentos''.
Quando a modificação é inevitável (por eficiência), documente
explicitamente esse contrato.
\end{boapratica}

\subsection{Aplicação prática: refinando o gerador de quebra-cabeças}

Vamos integrar este teste de unicidade ao gerador do Desafio~7 para
produzir Sudokus que tenham, por construção, solução única.

\begin{lstlisting}[caption={Gerador de Sudokus com solução única garantida.},label={lst:gerador_unico}]
/* Gera um quebra-cabeça com simetria reflexiva e solução única.    */
/* Retorna o número efetivo de células esvaziadas (pode ser menor    */
/* que o solicitado se o algoritmo esgotar tentativas).               */
int gerarSudokuValido(int s[DIM][DIM],
                      int completo[DIM][DIM],
                      int numAlvo)
{
    int i, j;
    int esvaziadas = 0;
    int tentativas = 0;
    const int MAX_TENT = 200;

    /* Começa com o tabuleiro completo. */
    for (i = 0; i < DIM; i++)
        for (j = 0; j < DIM; j++)
            s[i][j] = completo[i][j];

    while (esvaziadas < numAlvo && tentativas < MAX_TENT) {
        tentativas++;
        int ii = (rand() % LADO) + 1;
        int jj = (rand() % LADO) + 1;
        int i2 = (LADO + 1) - ii;
        int j2 = (LADO + 1) - jj;

        if (s[ii][jj] == VAZIO || s[i2][j2] == VAZIO)
            continue;

        /* Salva os valores antes de esvaziar. */
        int salvoA = s[ii][jj];
        int salvoB = s[i2][j2];

        s[ii][jj] = VAZIO;
        s[i2][j2] = VAZIO;

        /* Verifica unicidade. */
        if (contarUnicidade(s) == 1) {
            /* Pode manter o esvaziamento. */
            esvaziadas += (ii == i2 && jj == j2) ? 1 : 2;
            tentativas = 0;
        } else {
            /* Restaura para preservar unicidade. */
            s[ii][jj] = salvoA;
            s[i2][j2] = salvoB;
        }
    }
    return esvaziadas;
}
\end{lstlisting}

\begin{desempenho}
Este gerador é, deliberadamente, ``conservador'': prefere manter um
quebra-cabeça com solução única do que se arriscar a um Sudoku ambíguo.
O contador \codigo{tentativas} é zerado a cada esvaziamento bem-sucedido
e atinge \codigo{MAX\_TENT} apenas quando o algoritmo conclui que
remoções adicionais provocariam ambiguidade. O resultado típico contém
de $50$ a $60$ células preenchidas, faixa que abrange Sudokus de
dificuldade média a difícil.
\end{desempenho}

% =====================================================================
\section{Programa integrado e demonstração}

A Listagem~\ref{lst:demo} reúne os componentes principais em um único
\codigo{main} que demonstra o ciclo completo: geramos um Sudoku
completo via retrocesso randomizado, esvaziamos células simetricamente
preservando solução única e, ao final, resolvemos o quebra-cabeça
gerado.

\begin{lstlisting}[caption={Programa integrado.},label={lst:demo}]
/* Versão de resolverBacktracking que tenta dígitos em ordem        */
/* aleatória, produzindo Sudokus completos aleatórios.               */
int gerarCompletoAleatorio(int s[DIM][DIM])
{
    int i, j, k, d, tmp;
    int ordem[DIM];

    for (i = 1; i <= LADO; i++)
        for (j = 1; j <= LADO; j++)
            if (s[i][j] == VAZIO) {

                /* Constrói e embaralha a ordem dos dígitos 1..9. */
                for (k = 1; k <= LADO; k++) ordem[k] = k;
                for (k = LADO; k > 1; k--) {
                    int r = (rand() % k) + 1;
                    tmp = ordem[k]; ordem[k] = ordem[r]; ordem[r] = tmp;
                }

                for (k = 1; k <= LADO; k++) {
                    d = ordem[k];
                    if (ehMovimentoValido(s, i, j, d)) {
                        s[i][j] = d;
                        if (gerarCompletoAleatorio(s)) return 1;
                        s[i][j] = VAZIO;
                    }
                }
                return 0;
            }
    return 1;
}

int main(void)
{
    int completo[DIM][DIM] = {{0}};
    int puzzle[DIM][DIM];
    int copia[DIM][DIM];
    int i, j, n;

    srand((unsigned) time(NULL));

    /* 1) Gera um Sudoku completo aleatório. */
    if (!gerarCompletoAleatorio(completo)) {
        printf("Falha ao gerar Sudoku completo.\n");
        return 1;
    }
    printf("Sudoku completo gerado:");
    imprimirTabuleiro(completo);

    /* 2) Esvazia 50 células com simetria, preservando unicidade. */
    n = gerarSudokuValido(puzzle, completo, 50);
    printf("\nQuebra-cabeca com %d celulas esvaziadas:", n);
    imprimirTabuleiro(puzzle);

    /* 3) Resolve o quebra-cabeca para confirmar. */
    for (i = 0; i < DIM; i++)
        for (j = 0; j < DIM; j++) copia[i][j] = puzzle[i][j];

    if (resolverBacktracking(copia)) {
        printf("\nSolucao reconstruida:");
        imprimirTabuleiro(copia);
    } else {
        printf("\nERRO: o quebra-cabeca nao tem solucao.\n");
    }
    return 0;
}
\end{lstlisting}

\begin{boapratica}
Manter cópias separadas para ``solução de referência'' e ``cópia de
trabalho'' é a única forma robusta de verificar se um quebra-cabeça é
consistente. Tentar economizar memória reaproveitando o mesmo array
para os dois propósitos costuma se traduzir em horas de depuração.
\end{boapratica}

% =====================================================================
\section{Conclusão}

Neste apêndice, partimos da definição de um quebra-cabeça aparentemente
simples --- preencher uma grade de $9 \times 9$ com os dígitos de $1$
a $9$ obedecendo a três restrições --- e construímos, passo a passo,
soluções que vão da força bruta cega à programação de heurísticas
sofisticadas e à geração automática de novos quebra-cabeças. Mais que
um exercício de jogo, o Sudoku serviu como um campo de provas para
diversos paradigmas que serão utilizados ao longo de todo o livro:

\begin{itemize}
    \item \textbf{Encapsulamento via funções utilitárias.} Funções como
        \codigo{ehMovimentoValido} e \codigo{validSudoku} foram escritas
        \emph{uma única vez} e reaproveitadas em todos os solucionadores.
    \item \textbf{Recursão como ferramenta natural.} O retrocesso ficou
        elegantemente expresso por uma função recursiva curta.
        Voltaremos a esse padrão ao tratarmos do Passeio do Cavalo, das
        Oito Rainhas, da Torre de Hanói e da travessia de árvores.
    \item \textbf{Estruturas de iteração aninhadas.} O domínio bidimensional
        do Sudoku exigiu o uso disciplinado dos laços \codigo{for},
        \codigo{while} e \codigo{do...while} aprendidos no Capítulo~4.
    \item \textbf{Geração de números pseudoaleatórios.} A função \codigo{rand}
        permite explorar técnicas \textit{Monte Carlo} e introduzir variação
        nos geradores --- mas requer cuidado com a inicialização da semente.
    \item \textbf{Manipulação eficiente de conjuntos.} As máscaras de bits
        do Desafio~4 são um padrão idiomático em C, especialmente útil
        quando o universo é pequeno e fixo.
    \item \textbf{Heurísticas.} A heurística da Seção D.6 e a antecipação
        da Seção D.7 oferecem um primeiro contato com técnicas que dominam
        a inteligência artificial moderna.
    \item \textbf{Poda baseada no objetivo.} O Desafio~8 mostra como
        reformular uma pergunta cara em uma pergunta barata, decisão
        arquitetural que faz toda a diferença em problemas combinatórios.
\end{itemize}

\begin{engenharia}
A história da computação contém numerosos exemplos em que problemas
considerados ``insolúveis'' tornaram-se rotineiros graças não apenas
ao avanço do hardware, mas à invenção de algoritmos mais inteligentes.
O Sudoku é uma metáfora em miniatura desse processo: a força bruta
cega ($9^{81}$ possibilidades) jamais terminaria, ao passo que o
retrocesso bem implementado resolve qualquer Sudoku publicado em
milissegundos.
\end{engenharia}

Os recursos de programação aqui exercitados constituem, todos eles,
ferramentas com as quais o leitor já se tornou íntimo nos Capítulos~1
a~7. Cabe ao leitor agora estender este material: implementar duplas,
duplas ocultas, triplas e triplas ocultas; medir empiricamente quais
combinações de estratégias produzem os solucionadores mais eficientes;
gerar quebra-cabeças classificados por dificuldade; e --- por que não? ---
adicionar uma interface gráfica usando a biblioteca \textit{Allegro}
apresentada no Apêndice~E, que escapa do padrão da linguagem C mas
recompensa o trabalho com programas verdadeiramente bonitos.

Bom Sudoku!

\end{document}
